\subsection{Actual original homogeneous polarized weights at every normal-model point}

\begin{lemma}[Whole polynomial projection on all original model points]
\label{local:actual-full-point-polynomial-projection}
\lean{EnvelopeProperty.ProjSubvariety.map_normalizedSimultaneousPolynomialProjection_at}
\leanok
\uses{local:actual-normalized-full-point-restriction}
At EVERY true relative point $y$ of the SAME normalized simultaneous
model, its actual original polynomial-projective projection is exactly
the native valued point associated with the WHOLE normalized polynomial
point of its entire projected original product point.
\end{lemma}
\begin{proof}
Expand the literal original polynomial projection into normalization,
blowup, closed inclusion and inverse polynomial-product isomorphism.
The whole normalized/native homeomorphism inverse identity reduces the
other expression to exactly the same four native maps. The actual valued
point composition identity proves equality of entire points, including
the original residue valuations.
\end{proof}

\begin{definition}[True original homogeneous section norm on all model points]
\label{local:actual-full-point-homogeneous-section-norm}
\lean{EnvelopeProperty.ProjSubvariety.normalizedSimultaneousHomogeneousSectionNormAt, EnvelopeProperty.ProjSubvariety.normalizedSimultaneousHomogeneousSectionNormAt_eq_val}
\leanok
\uses{local:actual-full-point-polynomial-projection}
For each original homogeneous form $F$ of degree $d$, each original
coordinate-frame index $j$, and EVERY actual original model point $y$,
take the actual fiber section of the SAME globally constructed $O(d)$
section, pulled through the literal original polynomial projection.
Its TRUE model norm is EXACTLY $|F|_{\operatorname{res}y}$ on the SAME
original analytic variety restriction. This includes $F=0$ and all other
vanishing values; no nonvanishing hypothesis on $F$ is required.
\end{definition}
\begin{proof}
Apply the genuine proper model norm pullback theorem for the SAME actual
whole global section over the identity polynomial base map. The preceding
whole projection identity identifies the actual target point. The already
proved actual homogeneous-section norm theorem gives its entire
homogeneous polynomial value. The exact original coefficient-restriction
value formula identifies this with $|F|_{\operatorname{res}y}$. No metric
comparison is supplied as a hypothesis, and the original line and section
are retained through every actual pullback.
\end{proof}

\begin{definition}[Actual original polarized tensor section and weight]
\label{local:actual-full-point-polarized-ratio-section}
\lean{EnvelopeProperty.ProjSubvariety.normalizedSimultaneousPolynomialRatioSectionAt, EnvelopeProperty.ProjSubvariety.normalizedSimultaneousPolynomialRatioWeightAt}
\leanok
\uses{local:actual-full-point-homogeneous-section-norm, local:actual-normalized-full-point-ratio-section}
At every admissible actual original model point, tensor the preceding
TRUE original homogeneous fiber section with the SAME actual canonical
ratio section. This is a genuine generic fiber section of the SAME
original polynomial $O(d)$ tensor ratio line. Define its weight by its
TRUE negative-log model norm, using the already derived invertibility of
both actual original factors.
\end{definition}
\begin{proof}
Use the actual tensor product of the two genuine generic modules and its
canonical actual pure-tensor section. The original ratio and polynomial
line definitions identify the result with the SAME actual original
polarized line. Their actual invertibility proves invertibility of this
line; its actual model norm defines the stated weight.
\end{proof}

\begin{theorem}[Whole true polarized section-weight formula]
\label{local:actual-full-point-polarized-ratio-weight}
\lean{EnvelopeProperty.ProjSubvariety.normalizedSimultaneousPolynomialRatioWeightAt_eq}
\leanok
\uses{local:actual-full-point-polarized-ratio-section}
At EVERY such actual original model point and for EVERY original degree
$d$ form $F$, this SAME true polarized weight is the actual original
canonical ratio weight minus the ORIGINAL homogeneous $\log|F|$ at the
SAME original analytic restriction. The extended equality includes zero
sections and their infinite weights.
\end{theorem}
\begin{proof}
The true actual tensor model norm is the product of the two genuine
factor norms. Substitute the preceding original homogeneous section norm
formula. The actual extended logarithm sends this product of finite
nonnegative norms to the sum of the two extended logarithms. Neither
factor logarithm is positive infinity, so actual extended-real negation
and subtraction give exactly the displayed original ratio weight minus
original homogeneous logarithm. Zero factors are retained throughout.
\end{proof}
